\begin{abstract}
This work examines approaches to making computational models reversible. Broadly speaking, transforming a computational model into a reversible one, i.e. reversibilizing it, means extending its operational semantics conservatively in a way that each term of the model is interpretable as a bijection. 
We recall that the most common strategy to reversibilize a computational model yields operational semantics that halts computations whenever a computational state cannot be uniquely determined from its successor state, thereby allowing terms to be interpreted as \emph{partial} bijective functions.
We are interested in reversible computational models whose terms can be interpreted as \emph{total} bijective functions. This is essential for studying aspects of computational complexity related to reversible computational models.
We introduce \SCORE, a language designed for manipulating variables and stacks. 
Notably, common reversibilization strategies naturally lead to interpreting the functions for stack manipulation as partial bijections.
According to our interests, we demonstrate how to interpret \SCORE in a state space where, using a proof-assistant, we certify that stack operations are \emph{total} bijections. It follows that all \SCORE terms can be interpreted as \emph{total} bijections.
\end{abstract}